\chapter{实数完备性的等价形式}
\label{chap:equivalent-completeness}

\begin{chapterguide}
确界原理研究集合，单调收敛研究序列，区间套研究嵌套集合，Bolzano--Weierstrass 定理研究子列，Cauchy 完备性研究项与项之间的距离。本章证明这些命题在实数中可以相互推出，并标出 Archimedes 性质在“区间长度趋于零”及二分逼近中的作用。
\end{chapterguide}

本章临时使用\chapterref{chap:sequence-limit-definition}将系统展开的记号：\(a_n\to a\) 表示对每个 \(\varepsilon>0\)，从某项起 \(|a_n-a|<\varepsilon\)；子列写作 \(a_{n_k}\)，其中 \(n_1<n_2<\cdots\)。

\section{各种完备性命题及其依赖图}

考虑以下命题：
\begin{description}[style=nextline,leftmargin=3em]
 \item[(S) 确界原理] 每个非空有上界实数集有上确界。
 \item[(M) 单调有界收敛] 每个单调递增且有上界的实数列收敛；递减情形对称。
 \item[(N) 区间套] 若闭区间 \(I_n=[a_n,b_n]\) 嵌套且 \(b_n-a_n\to0\)，则交集恰含一个点。
 \item[(B) Bolzano--Weierstrass] 每个有界实数列都有收敛子列。
 \item[(C) Cauchy 完备] 每个实 Cauchy 列都收敛。
\end{description}

本章给出依赖链
\[
 \mathrm{(S)}\Longrightarrow\mathrm{(M)}
 \Longrightarrow\mathrm{(N)}
 \Longrightarrow\mathrm{(B)}
 \Longrightarrow\mathrm{(C)}
 \Longrightarrow\mathrm{(S)}.
\]
其中最后一个方向以及二分区间长度趋零的步骤显式使用 Archimedes 性质。

依赖图中的箭头表示一种无循环证明路线，不表示每个结论只能这样证明。例如可由确界原理直接证明 Bolzano--Weierstrass，也可从区间套经二分得到；可由 Cauchy 完备性与 Archimedes 网格直接重建单调收敛。比较不同路线时，必须把“有序域公理”“Archimedes 性质”和待证明的完备性命题分开记账。

\begin{proposition}[反例所说明的逻辑边界]
\(\Q\) 满足有序域公理和 Archimedes 性质，却不满足 (S)、(M)、(N)、(B)、(C) 的实数版本。因而 Archimedes 性质只负责尺度趋零，不能单独提供缺失极限。
\end{proposition}

\begin{proof}
\chapterref{chap:rational-defects}的有理 Cauchy 列直接否定 (C)。其单调有界十进制逼近列没有有理极限，否定 (M)；相应左右有理端点构成长度趋零而交集在 \(\Q\) 中为空的闭区间套，否定 (N)。无有理上确界的集合否定 (S)。若 (B) 在 \(\Q\) 中成立，该有理逼近列会有有理收敛子列，但平方误差迫使子列极限平方为 \(2\)，矛盾。
\end{proof}

\section{确界原理与单调有界收敛定理}

\begin{theorem}[单调有界收敛]
若 \((a_n)\) 单调递增且有上界，则
\[
 a_n\to\sup\{a_n:n\in\N\}.
\]
若单调递减且有下界，则它收敛到项集的下确界。
\end{theorem}

\begin{proof}
令 \(s=\sup\{a_n:n\in\N\}\)。给定 \(\varepsilon>0\)，\(s-\varepsilon\) 不是上界，所以存在 \(N\) 使 \(a_N>s-\varepsilon\)。单调性给出
\[
 s-\varepsilon<a_N\le a_n\le s\qquad(n\ge N),
\]
即 \(|a_n-s|<\varepsilon\)。递减情形对 \((-a_n)\) 应用递增结论。
\end{proof}

有界性用于取得有限确界；单调性把某一项进入误差区间的事实传递给整个尾部。删除任一条件结论都可能失败：\(a_n=n\) 单调但无界，\(a_n=(-1)^n\) 有界但不收敛。

\section{区间套定理及区间长度条件}

\begin{theorem}[区间套定理]
设 \(I_n=[a_n,b_n]\) 满足 \(I_{n+1}\subseteq I_n\)，且对每个 \(\varepsilon>0\)，从某项起 \(b_n-a_n<\varepsilon\)。则
\[
 \bigcap_{n=1}^{\infty}I_n=\{x\}
\]
恰含一个点。
\end{theorem}

\begin{proof}
左端点 \(a_n\) 单调递增且以上界 \(b_1\) 控制，故收敛到 \(x=\sup\{a_n\}\)。固定 \(n\)，每个 \(a_k\le b_n\)：当 \(k\le n\) 用 \(a_k\le a_n\le b_n\)，当 \(k\ge n\) 用 \(a_k\le b_k\le b_n\)。所以 \(x\le b_n\)，又 \(a_n\le x\)，故 \(x\in I_n\) 对所有 \(n\) 成立。

若 \(x,y\) 都在全部区间中，则 \(|x-y|\le b_n-a_n\) 对所有 \(n\) 成立。若 \(x\ne y\)，取 \(\varepsilon=|x-y|/2\)，长度条件给出矛盾，故唯一。
\end{proof}

若删除长度趋零条件，只能保证交集非空，不能保证单点。例如 \(I_n=[0,1]\) 的交仍为整个区间。若删除闭性，\((0,1/n)\) 的交为空。

\begin{proposition}[不要求长度趋零的闭区间套]
若非空闭区间 \(I_n=[a_n,b_n]\) 嵌套，则
\[
 \bigcap_{n=1}^{\infty}I_n
 =
 \left[\sup_n a_n,\inf_n b_n\right],
\]
右端可能退化为单点，但一定非空。
\end{proposition}

\begin{proof}
嵌套性使 \(a_n\) 递增、\(b_n\) 递减，并且每个 \(a_n\le b_m\)。故
\[
 a:=\sup_n a_n\le b:=\inf_n b_n.
\]
点 \(x\) 属于全部 \(I_n\) 当且仅当对每个 \(n\) 有 \(a_n\le x\le b_n\)，这等价于 \(a\le x\le b\)。
\end{proof}

\section{Bolzano--Weierstrass 定理}

\begin{theorem}[Bolzano--Weierstrass]
每个有界实数列都有收敛子列。
\end{theorem}

\begin{proof}
设 \(a_n\in[A,B]\)。把 \([A,B]\) 二分，至少一个半闭区间含有无穷多个序列项；选定这样的半区间为 \(I_1\)。继续二分，得到嵌套闭区间
\[
 I_k=[\alpha_k,\beta_k],\qquad
 \beta_k-\alpha_k=\frac{B-A}{2^k},
\]
每个都含无穷多个序列项。Archimedes 性质保证长度趋于零，区间套定理给出唯一交点 \(x\)。

递归选 \(n_k>n_{k-1}\) 且 \(a_{n_k}\in I_k\)，因为 \(I_k\) 含无穷多项。由于 \(x,a_{n_k}\in I_k\)，
\[
 |a_{n_k}-x|\le\beta_k-\alpha_k\longrightarrow0.
\]
故所得子列收敛到 \(x\)。
\end{proof}

“含无穷多个项”按指标计数，而非只看不同取值。常值数列的值集只有一个点，却有无穷多个指标可供抽取。

\begin{corollary}[有界无限集的聚点]
每个无限有界实数集 \(E\) 至少有一个聚点。
\end{corollary}

\begin{proof}
从 \(E\) 中递归选取互异点 \(x_n\)。该有界数列有收敛子列 \(x_{n_k}\to x\)。因子列各项互异，任意 \(x\) 的邻域都含异于 \(x\) 的 \(E\) 中点，故 \(x\) 是聚点。
\end{proof}

\section{Cauchy 完备性}

\begin{definition}
实数列 \((a_n)\) 称为 Cauchy 列，如果
\[
 \forall\varepsilon>0\ \exists N\in\N\
 \forall m,n\ge N,\quad |a_m-a_n|<\varepsilon.
\]
\end{definition}

\begin{proposition}
Cauchy 列有界。若 Cauchy 列有一个收敛到 \(a\) 的子列，则原列收敛到 \(a\)。
\end{proposition}

\begin{proof}
有界性证明与\chapterref{chap:cauchy-construction}相同。若 \(a_{n_k}\to a\)，给定 \(\varepsilon>0\)，先取 \(N\) 使 \(m,n\ge N\) 时 \(|a_m-a_n|<\varepsilon/2\)，再取 \(k\) 使 \(n_k\ge N\) 且 \(|a_{n_k}-a|<\varepsilon/2\)。对 \(n\ge N\)，
\[
 |a_n-a|\le|a_n-a_{n_k}|+|a_{n_k}-a|<\varepsilon.
\]
\end{proof}

\begin{theorem}[Cauchy 收敛准则]
实数列收敛当且仅当它是 Cauchy 列。
\end{theorem}

\begin{proof}
收敛蕴含 Cauchy 由三角不等式和 \(\varepsilon/2\) 得到。反之，Cauchy 列有界，由 Bolzano--Weierstrass 定理有收敛子列；上一命题把子列极限提升为原列极限。
\end{proof}

\begin{proposition}[尾部直径]
对数列 \((a_n)\) 定义
\[
 d_N=\sup\{\abs{a_m-a_n}:m,n\ge N\},
\]
允许 \(d_N=+\infty\)。则 \((a_n)\) 是 Cauchy 列当且仅当 \(d_N\) 最终有限且
\[
 d_N\to0.
\]
\end{proposition}

\begin{proof}
集合尾部随 \(N\) 增大而缩小，所以 \(d_N\) 单调不增。Cauchy 定义正是说：对每个 \(\varepsilon>0\)，存在 \(N\) 使该尾部中所有两项距离小于 \(\varepsilon\)，即 \(d_N\le\varepsilon\)（使用严格界时可先取 \(\varepsilon/2\)）。反向同理。
\end{proof}

尾部直径把双指标量词压缩为一个单调非负数列。对区间套端点，可取 \(d_N\le b_N-a_N\)；对快速误差算法，它直接给出停止准则。

\section{从 Cauchy 完备性返回确界原理}

\begin{theorem}
设 \(F\) 是 Archimedes 有序域。若 \(F\) 中每个 Cauchy 列都收敛，则 \(F\) 满足确界原理。
\end{theorem}

\begin{proof}
设 \(\varnothing\ne A\subseteq F\) 有上界。取 \(l_0\) 不是上界、\(u_0\) 是上界。已给 \(l_n,u_n\) 后令 \(m_n=(l_n+u_n)/2\)：若 \(m_n\) 是上界，取
\[
 l_{n+1}=l_n,\quad u_{n+1}=m_n;
\]
否则取
\[
 l_{n+1}=m_n,\quad u_{n+1}=u_n.
\]
始终有 \(l_n\) 不是上界、\(u_n\) 是上界，且
\[
 u_n-l_n=2^{-n}(u_0-l_0).
\]
Archimedes 性质使右端趋于零，所以两列都是 Cauchy 列。设 \(l_n\to l\)、\(u_n\to u\)。由 \(u_n-l_n\to0\) 得 \(l=u=:s\)。

若存在 \(a\in A\) 且 \(a>s\)，取 \(\varepsilon=(a-s)/2\)。充分大时 \(u_n<s+\varepsilon<a\)，与 \(u_n\) 是上界矛盾，故 \(s\) 是上界。若 \(c<s\)，充分大时 \(l_n>c\)；因 \(l_n\) 不是上界，存在 \(a\in A\) 使 \(a>l_n>c\)，所以 \(c\) 不是上界。故 \(s=\sup A\)。
\end{proof}

\section{Archimedes 条件在等价证明中的作用}

在实数中，Archimedes 性质已由确界原理推出，因此上述命题无须反复把它写入假设。反向从 Cauchy 完备性出发时，情况不同：二分宽度
\[
 2^{-n}(u_0-l_0)
\]
趋于零依赖自然数尺度能超过任意域元素。在非 Archimedes 有序域中，可能存在正无穷小 \(\eta\)，使 \(2^{-n}>\eta\) 对所有 \(n\) 成立；这时二分列未必逼近唯一边界。

因此严格的逻辑表述是：在 Archimedes 有序域中，(S)、(M)、(N)、(B)、(C) 等价；确界原理本身会推出 Archimedes 性质。只写“Cauchy 完备等价于确界完备”而不说明序与 Archimedes 背景，会掩盖条件。

\begin{remark}[序列完备与序完备]
Cauchy 完备只检验可数序列的距离稳定性，确界原理检验任意非空有界子集的序边界。在 Archimedes 有序域中，可用二分网格把集合边界编码成序列，二者因而等价；在非 Archimedes 环境中，这种可数网格未必能穿透无穷小尺度，所以两种完备性可以分离。
\end{remark}

\begin{proposition}
在 Archimedes 有序域中，只要每个单调递增有界数列收敛，就可推出确界原理。
\end{proposition}

\begin{proof}
对非空有上界集合 \(A\)，先取非上界 \(l_0\) 与上界 \(u_0\)，按正文二分构造 \(l_n,u_n\)。左端点单调递增有界，故收敛到 \(s\)；区间宽度趋零使右端点也趋于 \(s\)。正文最后两段只用“左端不是上界、右端是上界”即可证明 \(s=\sup A\)。
\end{proof}

\section*{本章小结}
\addcontentsline{toc}{section}{本章小结}

确界产生单调列极限；端点极限给出区间套；二分区间抽取收敛子列；收敛子列使 Cauchy 列整体收敛；最后以“非上界—上界”二分列重建上确界。长度趋零把序逼近与距离逼近连接起来，其依据是 Archimedes 性质。

\section*{习题}
\addcontentsline{toc}{section}{习题}

\noindent\textbf{配置说明：}本章按I级枢纽章配置，共26道编号题，含49个独立训练单元，建议总用时10至14小时。

\subsection*{A组　各形式的直接应用}
\begin{enumerate}[label=\textbf{6.\arabic*.}]
 \exerciseitem{6.1} \mustdo 从定义证明收敛数列是 Cauchy 列。
 \exerciseitem{6.2} \mustdo 证明单调递增列若无上界，则对每个 \(M\in\R\) 从某项起 \(a_n>M\)。
 \exerciseitem{6.3} \mustdo 给出有界但不收敛数列、单调但不收敛数列，并指出各自缺少的条件。
 \exerciseitem{6.4} \mustdo 证明区间套定理中若只知交集非空而无长度条件，交集必为闭区间。
 \exerciseitem{6.5} \mustdo 分别给出删除嵌套性、闭性、长度趋零条件后结论失败的例子。
\end{enumerate}

\subsection*{B组　子列与 Cauchy 性}
\begin{enumerate}[label=\textbf{6.\arabic*.},resume]
 \exerciseitem{6.6} \mustdo 证明有界数列若只有一个可能的子列极限，则原列收敛。
 \exerciseitem{6.7} \mustdo 证明 Cauchy 列的任意子列仍是 Cauchy 列；若某子列收敛，则所有收敛子列极限相同。
 \exerciseitem{6.8} \mustdo 用二分区间法证明每个无限有界实数集至少有一个聚点。
 \exerciseitem{6.9} \mustdo 证明若 \(|a_{n+1}-a_n|\le2^{-n}\)，则 \((a_n)\) 是 Cauchy 列。
 \exerciseitem{6.10} \optionaldo 构造 \(|a_{n+1}-a_n|\to0\) 但不是 Cauchy 列的数列。
 \exerciseitem{6.11} \mustdo 证明有限多个 Cauchy 列的线性组合仍为 Cauchy 列。
\end{enumerate}

\subsection*{C组　等价性项目}
\begin{enumerate}[label=\textbf{6.\arabic*.},resume]
 \exerciseitem{6.12} \projectdo\ 【前置：本章全部主定理；预计用时：90分钟】不经过 Bolzano--Weierstrass 定理，直接用快速 Cauchy 子列与级数型误差证明 Cauchy 完备性。
 \exerciseitem{6.13} \mustdo 补全从 Cauchy 完备性构造上确界时 \(l_n,u_n\) 的 Cauchy 性和共同极限证明。
 \exerciseitem{6.14} \mustdo 证明“每个有界数列都有 Cauchy 子列”与 Bolzano--Weierstrass 定理在 Cauchy 完备的有序域中等价。
 \exerciseitem{6.15} \optionaldo 证明闭区间的可数开覆盖若具有 Lebesgue 数，则可抽取有限子覆盖；说明区间套思想的作用。
 \exerciseitem{6.16} \opendo 调查一个非 Archimedes 有序域，说明 Cauchy 完备、序完备与球完备三种概念的区别。
 \exerciseitem{6.17} \mustdo 证明在 \(\Q\) 中 (S)、(M)、(N)、(B)、(C) 均失败，并尽量使用\chapterref{chap:rational-defects}同一个 \(\sqrt2\) 缺口。
 \exerciseitem{6.18} \mustdo 对不要求长度趋零的嵌套闭区间证明
 \[
  \bigcap_n[a_n,b_n]=[\sup_n a_n,\inf_n b_n].
 \]
 \exerciseitem{6.19} \mustdo 从 Bolzano--Weierstrass 定理推出每个无限有界实数集有聚点，并解释互异点选择的必要性。
 \exerciseitem{6.20} \mustdo 定义尾部直径 \(d_N\)，证明 Cauchy 性等价于 \(d_N\to0\)。
 \exerciseitem{6.21} \mustdo 证明若 \((a_n)\) 是 Cauchy 列且有子列趋于 \(+\infty\)，则产生矛盾；说明 Cauchy 列不能有不同的扩充实子列极限。
 \exerciseitem{6.22} \mustdo 只用单调有界收敛与 Archimedes 性质重建确界原理。
 \exerciseitem{6.23} \mustdo 证明若每个长度趋零的闭区间套交集非空，则交集自动为单点；据此区分“存在性”与“唯一性”的条件来源。
 \exerciseitem{6.24} \optionaldo 给出从 (C) 直接推出 (B) 的证明方案，避免先使用确界原理；可先构造 Cauchy 子列。
 \exerciseitem{6.25} \projectdo\ 【前置：本章全部内容；预计用时：120分钟】为依赖图中每条箭头写一份条件清单和反例清单，确保不循环引用等价命题。
 \exerciseitem{6.26} \opendo 比较可数区间套与任意全序嵌套闭区间族的交性质，说明为何后者会导向比序列完备更强的完备概念。
\end{enumerate}

\section*{部分习题提示}
\begin{itemize}
 \item \exerciseref{6.6}：反设不收敛，抽取一条始终远离候选极限的子列，再应用收敛子列定理。
 \item \exerciseref{6.10}：取调和部分和 \(a_n=\sum_{k=1}^n1/k\)。
 \item \exerciseref{6.12}：抽取 \(n_k\) 使尾部项两两距离 \(<2^{-k}\)，证明它收敛后再控制原列。
\end{itemize}
